% Authors: Keerthi Rapolu & Sreeja Katta

IFS maps the embedding cosine to four alignment categories: \emph{well-aligned}
($\mathrm{IFS} \geq 0.85$), \emph{minor divergence} ($0.70 \leq \mathrm{IFS}
< 0.85$), \emph{significant divergence} ($0.50 \leq \mathrm{IFS} < 0.70$), and
\emph{severe divergence} ($\mathrm{IFS} < 0.50$). These thresholds are chosen
so that the IBD flag ($\mathrm{IFS} < 0.65$) falls between the minor and
significant bands, giving anomaly detection a conservative operating point that
does not require separate threshold calibration.

The score decomposes into four interpretable sub-scores — type alignment,
utilization alignment, duration alignment, and resource alignment — weighted
0.35, 0.25, 0.20, and 0.20 respectively. For LLM pipeline workloads, a
token-waste sub-score replaces part of the resource term because token budget
adherence carries governance significance that node count alone does not capture.
These sub-scores are presented in the dashboard as interpretability support
rather than as standalone metrics; the canonical IFS value remains the embedding
cosine.

The role of IFS within PBCP is diagnostic rather than preventive. A workload can
receive high CPS because an effective intervention was applied and high IFS
because runtime behavior remained aligned, or low CPS because no intervention
was triggered and low IFS because the job behaved contrary to its stated intent.
That separation is what makes IFS useful: it explains why waste occurred rather
than only measuring whether it was stopped, and it gives the policy learning loop
a signal that CPS and ESR do not directly capture.
